\documentclass[10pt]{article}

\usepackage[utf8]{inputenc}

\usepackage[T1]{fontenc}

\usepackage{amsmath}

\usepackage{amsfonts}

\usepackage{amssymb}

\usepackage[version=4]{mhchem}

\usepackage{stmaryrd}

\usepackage{bbold}



\begin{document}

ленды е и (алге б р аически е) раз ве тв ленны е әлем енты.



Допуская гри продолжении исходного элемента $\left(U(a, R), f_{a}\right)$ и особые әлементы с рядами вида (3), вообце говоря, многозначные (при $v>1$ ) и имеющие особенности типа полюса (при $\mu<0$ ), получают более обширную, чем вейерштрассова, р и м а н о в у о бла с т ь с у ш е с т в о в а ни я $E_{R}$ и соответствующую более обширную совокупность элементов, определяемых рядами вида (3), называемую аналитическим образом. Аналитич. образ отличается от П. а. ф. присоединением всех особых әлементов, получаемых при продолжении данного регулярного элемента. После введения соответствующей топологии аналитич. образ превращается в риманову поверхность данной функции.



При описанном построении П. а. ф. $f_{W}$ можно пользоваться вместо элемента понятием ростка аналитич. функции, смысл введения к-рого заключается в локализации понятия әлемента, в отвлечении от не имеющей в данном случае существенного значения величины радиуса сходимости. Два элемента $(D, f)$ и $(G, h)$ такие, что области $D$ и $G$ содержат общую точку $a$, наз. э к в и в а л е н т ны м и в то чк е $a$, если существует окрестность точки $a$, в к-рой $f \equiv h$. Это отношение эквивалентности обладает обычными свойствами рефлексивности, симметрии и транзитивности. Класс эквивалентности элементов в данной точке $a \in \overline{\mathbb{C}}$ наз. рост ком ан а литичес ко й функц и и $\mathbf{f}_{a}$ в точке $a$. Росток характеризует локальныге свойства функции в данной точке. Два ростка $f_{a}$ и $\mathbf{g}_{a}$ равны, если в нек-рой окрестности точки $a$ совпадают какие-либо представители классов эквивалентности. Аналогично, при помощи представителей, определяются арифметич. действия с ростками и их дифференцирование. П. а. Ф. $f W$ есть совокупность всех ростков аналитич. функций $\mathbf{f}_{\zeta}, \zeta \in E_{f}$, получаемых из данного ростка $\mathbf{f}_{a}$ аналитич. продолжением вдоль всевозможных путей в $\overline{\mathbb{C}}$. Равенство двух П. а. ф. $f_{W}, g_{W}$ и действия с П. а.ф. определяются как равенство ростков $\mathbf{f}_{a}, \mathbf{g}_{a}$ в какой-либо точке $a \in E_{f} \cap E_{g}$ и действия с ростками.



Элементы $(D, f)$, вейерштрассовы элементы ( $U^{n}(a$, $R), f_{a}$ ) и ростки аналитич. функции многих комплексных переменных $z=\left(z_{1}, \ldots, z_{n}\right), n \geqslant 1$, определяются точно так же, как выше, но с помощью областей $D$ комплексного пространства $\mathbb{C}^{n}$ или поликругов сходимости



$U^{n}(a, R)=\left\{z \in \mathbb{C}^{n}:\left|z_{j}-a_{j}\right|<R_{j}, j=1, \ldots, n\right\} ;$



$R_{j}>0, j=1, \ldots, n ; a=\left(a_{1}, \ldots, a_{n}\right)$; $R=\left(R_{1}, \ldots, R_{n}\right)$,



кратных степенных рядов



$f_{a}=f_{a}(z)=\sum_{k_{1}, \ldots, k_{n}=0}^{\infty} c_{k_{1} \ldots k_{n}}\left(z_{1}-a_{1}\right)^{k_{1}} \ldots\left(z_{n}-a_{n}\right)^{k}$.



Понятие П. а. ф. многих комплексных переменных строится далее вполне аналогично случаю одного переменного.



Лum.: [1] М а р к уш е в и ч А. И., Теория аналитических функций, 2 изд., т. 2, М., 1968; [2] III а б а т Б. В., Введение в комплексный анализ, 2 изд., ч. 1-2, М., 1976; [3] С п p и нг е р Д ж., Введение в теорию римановых поверхностей, пер. с англ., М., 1960; 4] Ф у к с Б. А., Введение в теорию анали-



Е. Д. Соломенцев.



ПОЛНАЯ ВАРИАЦИЯ Ф у к ц и и - то же, что вариация бункции.



ПОЛНАЯ ГРУППА - группа, в к-рой для любого ее элемента $g$ и любого целого числа $n \neq 0$ разрешимо уравнение $x^{n}=g$. Абелева П. г. наз. также д е лим о й г р у п п о й. Важными примерами П. г. являются аддитивная группа всех рациональных чисел и группа всех комплексных корней из 1 степеней $p^{k}, k=1$, $2, . .$, где $p$ - простое чісло (к в а з и ци и л и ч ес к а я г р у п п а). Всякая абелева П. г. разлагается в прямую сумму груюп, каждая из к-рых изоморфна одной из указанных. О неабелевых П. Г. известно значительно меньше. Всякая неединичная П. г. бесконечна. Всякая группа вложима в подходящую П. г. Если в П. г. указанные в определении уравнения разрешимы однозначно, она наз. $D-\Gamma$ р у п п о й. Таковы, в частности, локально нильпотентные П. г. без кручения.



[2] Л. К а $:$ [1] К у р о ш А. Г., Теория групп, 3 изд., М., 1967; теории групп, 3 изд., М. 1982 . А. Л. Нмелькин.



ПОЛНАЯ КРИВИЗНА - 1) П. к. в т о ч к е пов е р х н о т и $\Phi$ в евклидовом пространстве $\mathbb{R}^{3}-$ скалярная величина $K$, равная произведению главных (нормальных) кривизн $k_{1}$ и $k_{2}$, вычисляемых в точке поверхности: $K=k_{1} k_{2}$; наз. также гауссовой кривизной поверхности. Понятие П. к. обобщается для гиперповерхности в евклидовом пространстве $\mathbb{R} n+1, n>2$. П. к. в этом случае есть величина $K=k_{1}$. . $k_{n}$, где $k_{i}-$ главная нормальная кривизна в точке гишерповерхности в $i$-м главном направлении.



П. к. в точке двумерной поверхности в трехмерном римановом пространстве равка разности внутренней кривизны - римановой кривизны двумерной поверхности, и внешней кривизны - римановой кривизны объемлющего пространства в направлении бивектора, касательного к поверхности в рассматриваемой точке.



\begin{enumerate}

  \setcounter{enumi}{1}

  \item П. к. о б л а с т и $D$ на поверхности Ф в евклидовом пространстве $\mathbb{R}^{3}$ - величина $\iint_{D} K d \sigma$, где $K-$ гауссова кривизна поверхности в точке, $d \sigma-$ элемент площади поверхности. Аналогично определяется П. к. области нек-рого риманова многообразия, причем под $K$ понимается риманова кривизна многообразия, вычисляемая в точках многообразия в направлении касательных бивекторов, а интегрирование ведется по площади (мере) области многообразия. Л. А. Сидоров.

\end{enumerate}



ПОЛНАЯ ЛИНЕИНАЯ ГРУПІА - группа всех обратимых матриц степени $n$ над ассоциативным кольцом $K$ с единицей; общепринятое обозначение: $\mathrm{GL}_{n}(K)$ или $\mathrm{GL}(n, K)$. П. л. г. $\mathrm{GL}(n, K)$ может быть также определена как группа автоморфизмов $\operatorname{Aut}_{K}(V)$ свободного правого $K$-модуля $V$ с $n$ образугщими.



В исследовании группы $\mathrm{GL}(n, K)$ большой интерес представляет вопрос о ее нормальном строении. Центр $Z_{n}$ группы $\mathrm{GL}(n, K)$ состоит из скалярных матриц с элементами из центра кольца $K$. В классич. случае, когда $K$ - поле, решающую роль играет исследование нормального строения специальной линейной группы $\operatorname{SL}(n, K)$, состоящей из матриц с определителем 1. А именно, коммутант группы $\mathrm{GL}(n, K)$ совпадает с $\mathrm{SL}(n, K)$ (кроме случая $n=2,|K|=2)$, и всякая нормальная подгруппа группы $\mathrm{GL}(n, K)$ либо содегжится в $Z_{n}$, либо содержит $\operatorname{SL}(n, K)$. В частности, специальная проективная г ру пи п $\operatorname{PSL}(n, K)=\operatorname{SL}(n, K) / \mathrm{SL}(n, K) \cap Z_{n}$.



является простой (за исключением случаев $n=2$, $|K|=2,3)$.



Если $K$ - тело и $n>1$, то всякая нормальная подгруппа группы $\mathrm{GL}(n, K)$ либо содержится в $Z_{n}$, либо содержит коммутант $\mathrm{SL}+(n, K)$ группы $\mathrm{GL}(n, k)$, причем коммутант $\mathrm{SL}+(n, K)$ порождается трансвекциями и факторгруппа $\mathrm{SL}+(n, K) / \mathrm{SL}+(n, K) \cap Z_{n}$ проста. Кроме того, существует естественный изоморфизм



\[

\mathrm{GL}(n, K) / \mathrm{SL}^{+}(n, K) \simeq K^{*} /\left[K^{*}, K^{*}\right],

\]



где $K^{*}-$ мультипликативная группа тела $K$. Если $K$ конечномерно над своим центром $k$, то роль группы





\end{document}